\subsubsection*{Resultados de la Regresión Logística}

Tras la evaluación con el conjunto de prueba (20 muestras), el modelo alcanzó un rendimiento perfecto:
\begin{itemize}
    \item \textbf{Precisión Global (Accuracy):} 100\%.
    \item \textbf{Reporte de Clasificación:} Tanto para la clase \textit{Fake} como \textit{Real}, se obtuvieron valores de 1.00 en precisión, \textit{recall} y F1-score.
    \item \textbf{Matriz de Confusión:} Se clasificaron correctamente las 10 muestras de cada clase sin falsos positivos ni falsos negativos.
\end{itemize}
\textbf{Conclusión:} La Regresión Logística demuestra que el conjunto de características acústicas seleccionado es altamente discriminativo para este problema, permitiendo una separación lineal perfecta de las clases en el entorno de prueba. Por otro lado es extraño que den resultados tan perfectos, pero al escuchar algunos de los audios generados descubrimos que son muy dicriminatorios y por tanto los puede identificar fácilmente.
